% !TeX root = ../main.tex
\section{结论}
本文按 V2、V3、V4 的顺序研究 BG6022 在自主量子化学计算中的三种思路。V2 通过受控工作流建立可复核计算链，V3 在有限任务与对话情境中实现语义参与、工具组合和有界修复，V4 将现场生成代码纳入证据约束闭环，使任务表达从预定义工具扩展至程序化算法。

V4 是本研究的重点。单决策循环、目标与计划分离、统一 ORCA 执行以及结果资格分层形成了受控开放任务的方法基础。现有证据在 T1--T5 与 H1 六类任务中保存了成功 ORCA 计算与实际数值，包括乙醇的无虚频优化结构、正己烷有限候选能量、正丁烷六个约束优化点、丙酮五态激发、两种基组下的环境电子能差和 HCN 屏蔽张量。原冻结批次 8 次启动中 7 次计算成功，成功比例为 87.5\%；H2 尚无极化率数值。这些结果支持代码生成路线的跨任务执行能力，后续研究应进一步完善几何核查、结果提取、完整交付与重复试验。

\paragraph{数据与材料可用性}
支撑本文图表的结构化数据、来源校验信息、绘图程序与 LaTeX 源文件作为补充材料提供。本文对既有实验作回顾性分析，本次文档修订未启动新的计算实验，补充数值来自已保存的后续开发批次。部分早期记录的完整执行快照与独立数值参考不可获得，相应限制已在实验评价中说明。架构示意图由 GPT Image 辅助绘制，六幅量化图由记录数据直接计算生成。
